\section{Illustration: Same Expertise, Different Decision Value}
\label{sec:illustration}

To make the divergence between predictive expertise and decision value concrete, consider a binary classification problem ($K = 2$, $|\mathcal{A}| = 2$: discharge or admit) under two reward structures.

\paragraph{Setup.}
The model produces belief $b_x = P(Y\!=\!1 \mid x)$.
A human reviewer observes the instance and shifts the belief to $b_{x,h}$.
We compare two reward matrices:
\begin{itemize}[nosep]
  \item \textbf{R1} (symmetric accuracy): $R = \bigl[\begin{smallmatrix}1&0\\0&1\end{smallmatrix}\bigr]$, threshold $b^* = 0.50$.
  \item \textbf{R2} (FP penalty): $R = \bigl[\begin{smallmatrix}1&0\\{-4}&1\end{smallmatrix}\bigr]$, threshold $b^* = \tfrac{5}{6} \approx 0.833$.
\end{itemize}

\paragraph{Two cases.}
\textbf{Case~A} (high expertise, low decision value): $b_x = 0.55$, $b_{x,h} = 0.70$.
The physician shifts the belief by $+0.15$---a substantial probabilistic update.

\textbf{Case~B} (low expertise, high decision value): $b_x = 0.80$, $b_{x,h} = 0.87$.
The physician shifts the belief by only $+0.07$---half the absolute update of Case~A.

A residual-expertise audit (log-loss gain at $y\!=\!1$: $0.24$ for Case~A vs.\ $0.08$ for Case~B) ranks Case~A as the far more valuable human signal.

\paragraph{Boundary regret tells a different story.}
Under \textbf{R1} (threshold $0.50$): both cases have the model action ``admit'' ($b_x > 0.5$), and neither human update changes the action.
$\BR = 0$ for both.

Under \textbf{R2} (threshold $0.833$): the model discharges in both cases ($b_x < 0.833$).
In Case~A, $b_{x,h} = 0.70 < 0.833$: the human-updated belief remains below the threshold and $\BR = 0$.
In Case~B, $b_{x,h} = 0.87 > 0.833$: the posterior crosses the decision boundary.
Computing:
\[
  \BR_{\text{B}} = V(0.87) - r_{\text{discharge}} \cdot (1 - 0.87, \; 0.87)
  = 0.35 - 0.13 = 0.22.
\]

\noindent
Case~A, with three times the log-loss improvement, has zero decision value under R2.
Case~B, with modest predictive improvement, carries all the decision value---because its update is located near the reward-induced boundary.

\paragraph{Takeaway.}
A deployment-focused audit must measure not how much the human improves the probability estimate, but whether that improvement is located near a decision facet.
The same human signal can be high-expertise/zero-decision-value or modest-expertise/high-decision-value, depending entirely on the reward geometry.
This is the distinction that boundary regret formalizes, and that Experiments~2 and~3 (Sections~\ref{sec:exp2} and~\ref{sec:exp3}) demonstrate at scale.
